\section{项目特色与预期改进}

\subsection{项目特色}

\begin{enumerate}
    \item \textbf{清晰的数据到决策分析主链：}本课题不堆砌方法，而是以一条清晰主线组织工作——``跨年度数据审计 $\rightarrow$ 能耗基准预测 $\rightarrow$ 异常高耗能识别 $\rightarrow$ 建筑级汇总 $\rightarrow$ 改造优先级排序 $\rightarrow$ 辅助聚类画像''。每个环节有明确的输入输出和质量检查节点，聚类降级为辅助画像而非核心目标。区别于大多数仅训练单一模型或简单排列方法的课题。

    \item \textbf{按年份滚动验证与双重泛化检验：}主实验采用按完整年份的扩展窗口滚动验证（2015--2017$\rightarrow$2018至2015--2021$\rightarrow$2022），更真实地模拟``用历史数据估计基准''的场景；辅助实验采用按OSEBuildingID的GroupShuffleSplit，检验模型对从未见过的新建筑的泛化能力。两种验证策略互补，全面评估模型性能。

    \item \textbf{三层递进式异常识别机制：}不再使用``残差$>0$''直接判定异常，而是采用三层递进筛选——（1）模型预测残差异常（同类型同年度标准化残差$>1.5$或前10\%）；（2）同类P75基准超标；（3）2022--2024年持续性验证。三层交叉验证可显著降低单一方法的误判率。

    \item \textbf{建筑级记录聚合与持续性评估：}原始数据为``建筑—年份''粒度，本研究按OSEBuildingID将多年记录聚合为建筑级结果（2024年当前状态 + 2022--2024年持续性表现），每栋建筑仅保留一行。持续高耗能频率和最近三年平均残差等指标使优先级排序更加稳健。

    \item \textbf{可配置的优先级评分与全面的稳定性检验：}设计6维度加权评分框架（标准化残差25\%、基准超耗20\%、持续频率15\%、超额能耗量15\%、GHG强度15\%、建筑年龄10\%），引入数据质量折扣因子防止缺失严重的建筑因极端值错误排前；通过Bootstrap、Jaccard相似度、Spearman/Kendall秩相关系数等多维度检验排序稳定性。支持4套权重方案对比（标准/节能优先/减碳优先/大型建筑优先）。

    \item \textbf{全面的评估基线与分组报告：}以全体均值、建筑类型中位数作为基线，当且仅当复杂模型显著优于简单基线时才采信其结果；评估时按年份、建筑类型、高能耗子集、已见/未见建筑分组报告性能指标。

    \item \textbf{高度模块化的代码结构：}7个独立Python脚本，每个模块职责单一、输出明确，可独立运行或通过主控脚本一键执行。所有预处理拟合严格限定在训练集，避免数据泄漏。

    \item \textbf{完善的跨年度数据审计：}在正式建模前对各年份数据进行系统性审计（字段名称一致性、单位一致性、建筑类型命名一致性、ID稳定性、重复记录、数值解析、字段覆盖率），生成字段映射表和覆盖率统计等审计文件，确保数据质量可追溯。
\end{enumerate}

\subsection{与现有工作的区别}

\begin{enumerate}
    \item 现有建筑能耗benchmarking项目多使用ENERGYSTAR Score作为主要依赖（该分数本身由当年能耗数据计算），本研究设置两组实验——主模型（模型A）不使用ENERGYSTAR Score，仅保留ES\_Missing指示变量，模型B作为扩展对比，更真实地模拟``未获评建筑''的预测场景并防止模型性能因评级信息而虚高；

    \item 将基准性能范围数据（Benchmarking Performance Ranges by Building Type）作为异常识别第二层（同类基准超标）的参照系，而非仅用于EDA展示；

    \item 在优先级排序中采用三层递进筛选（预测残差异常 + 同类基准超标 + 持续性验证），而非简单的正残差截断；同时引入数据质量折扣因子，避免因单一极端值导致的错误推荐；

    \item 不同建筑类型天然能耗水平差异巨大（如医院vs仓库），直接用绝对SiteEUI排序会导致某些类型垄断Top-50，通过同类型超耗比例替代绝对能耗，使排序更加公平；

    \item 将聚类从``识别异常''的核心角色降级为``候选建筑画像''的辅助角色，聚类对象仅为高优先级候选建筑，使用连续数值特征计算距离，建筑类型仅用于事后解释，逻辑更加清晰。
\end{enumerate}

\subsection{待改进方面}

\begin{enumerate}
    \item \textbf{数据扩充可能性：}西雅图公开数据未包含墙体保温、窗户类型等工程细节，若引入EnergyPlus模拟或结合气候数据（如采暖/制冷度日数），可显著提升预测精度。

    \item \textbf{评分权重可进一步优化：}当前权重为手动设定并通过敏感性分析验证，如引入实际改造成本数据（如ECM成本面板），可将主观排序目标转化为成本效益优化问题。

    \item \textbf{因果关系需慎重声明：}SHAP显示的特征重要性属于模型内部关联，不代表真实因果关系；在实际政策决策中需结合领域知识谨慎解读。报告中统一使用``重要预测特征''表述。

    \item \textbf{聚类稳定性：}K-Means对初始中心敏感且假设各向同性，如使用DBSCAN或谱聚类可能发现非凸形状的聚类结构。但由于聚类已降级为辅助画像，该问题对核心结论的影响有限。

    \item \textbf{滞后特征未纳入：}当前主模型仅使用静态建筑属性和年度背景变量，未加入上一年SiteEUI、近三年均值、同比变化率等滞后特征。若后续研究需要严格的时间序列预测，可在此基础上扩展。
\end{enumerate}
